\documentclass[12pt, a4paper]{article}

% ── Packages ──────────────────────────────────────────────────────────────────
\usepackage[utf8]{inputenc}
\usepackage[T1]{fontenc}
\usepackage{lmodern}
\usepackage[margin=2.5cm]{geometry}
\usepackage{xcolor}
\usepackage{hyperref}
\usepackage{enumitem}
\usepackage{tabularx}
\usepackage{booktabs}
\usepackage{tcolorbox}
\usepackage{titlesec}
\usepackage{fancyhdr}
\usepackage{parskip}
\usepackage{microtype}
\usepackage{amssymb}

% ── Colors ────────────────────────────────────────────────────────────────────
\definecolor{emerald}{HTML}{10B981}
\definecolor{emeralddark}{HTML}{059669}
\definecolor{slate}{HTML}{334155}
\definecolor{slatelight}{HTML}{64748B}
\definecolor{bluehint}{HTML}{3B82F6}
\definecolor{amberbg}{HTML}{FFF7ED}
\definecolor{amberframe}{HTML}{F59E0B}
\definecolor{redflag}{HTML}{EF4444}
\definecolor{greenbg}{HTML}{F0FDF4}
\definecolor{greenframe}{HTML}{22C55E}
\definecolor{bluebg}{HTML}{EFF6FF}
\definecolor{blueframe}{HTML}{3B82F6}

% ── Hyperref Setup ────────────────────────────────────────────────────────────
\hypersetup{
  colorlinks=true,
  linkcolor=emeralddark,
  urlcolor=bluehint,
  pdftitle={Interview Questions — Expense Tracker Project},
  pdfauthor={Senior Interviewer},
}

% ── Section Styling ───────────────────────────────────────────────────────────
\titleformat{\section}
  {\Large\bfseries\color{emeralddark}}
  {Round \thesection}{1em}{}
  [\vspace{-0.5em}\textcolor{emerald}{\rule{\textwidth}{1.5pt}}]

\titleformat{\subsection}
  {\large\bfseries\color{slate}}
  {Q\thesubsection.}{0.6em}{}

% ── Header / Footer ──────────────────────────────────────────────────────────
\pagestyle{fancy}
\fancyhf{}
\fancyhead[L]{\small\textcolor{slatelight}{Interview Questions — Expense Tracker}}
\fancyhead[R]{\small\textcolor{slatelight}{Fresher / Junior Candidate}}
\fancyfoot[C]{\small\textcolor{slatelight}{\thepage}}
\renewcommand{\headrulewidth}{0.4pt}
\renewcommand{\footrulewidth}{0pt}

% ── Custom Environments ──────────────────────────────────────────────────────
\newtcolorbox{listenfor}{
  colback=greenbg, colframe=greenframe,
  coltitle=white, fonttitle=\bfseries,
  title={What I'm Listening For},
  boxrule=0.8pt, arc=4pt, left=8pt, right=8pt, top=6pt, bottom=6pt,
}

\newtcolorbox{followup}{
  colback=bluebg, colframe=blueframe,
  coltitle=white, fonttitle=\bfseries,
  title={Follow-Up Probe},
  boxrule=0.8pt, arc=4pt, left=8pt, right=8pt, top=6pt, bottom=6pt,
}

\newtcolorbox{tipbox}{
  colback=amberbg, colframe=amberframe,
  coltitle=white, fonttitle=\bfseries,
  title={Interviewer Tip},
  boxrule=0.8pt, arc=4pt, left=8pt, right=8pt, top=6pt, bottom=6pt,
}

% ══════════════════════════════════════════════════════════════════════════════
\begin{document}

% ── Title Page ────────────────────────────────────────────────────────────────
\begin{titlepage}
  \centering
  \vspace*{4cm}

  {\Huge\bfseries\textcolor{emeralddark}{Interview Questions}}\\[0.6cm]
  {\LARGE\textcolor{slate}{Expense Tracker — MERN Project}}\\[1.8cm]

  {\large\textcolor{slatelight}{For a fresher / candidate with very little experience}}\\[3cm]

  \textcolor{emerald}{\rule{8cm}{1.5pt}}\\[1.2cm]

  {\normalsize\textcolor{slatelight}{%
    \textbf{Interviewer persona:} Senior full-stack engineer, 10 years of experience.\\[0.4cm]
    \textbf{Philosophy:} I'm not trying to trick or stump the candidate.\\
    I want to understand if they \emph{actually built} this, how deeply they\\
    understand what they wrote, and whether they can reason about\\
    trade-offs — even if they don't always have the ``right'' answer.%
  }}

  \vfill
  {\small\textcolor{slatelight}{Total Duration: $\sim$40--50 minutes \quad|\quad 6 Rounds \quad|\quad 20 Questions}}
\end{titlepage}

% ── Table of Contents ─────────────────────────────────────────────────────────
\tableofcontents
\newpage

% ── Overview Table ────────────────────────────────────────────────────────────
\begin{center}
\renewcommand{\arraystretch}{1.4}
\begin{tabularx}{\textwidth}{>{\bfseries}l X c c}
  \toprule
  \textbf{Round} & \textbf{Focus Area} & \textbf{Duration} & \textbf{Questions} \\
  \midrule
  1 & Warm-Up \& Project Understanding   & 5--8 min   & 3 \\
  2 & Backend \& Database Deep-Dive       & 10--12 min & 5 \\
  3 & Frontend \& React                   & 8--10 min  & 4 \\
  4 & API Design \& Error Handling        & 5--7 min   & 3 \\
  5 & DevOps \& Deployment                & 5--7 min   & 3 \\
  6 & Thinking \& Growth                  & 5 min      & 2 \\
  \midrule
  & \textbf{Total} & \textbf{$\sim$40--50 min} & \textbf{20} \\
  \bottomrule
\end{tabularx}
\end{center}

\vspace{1cm}

% ══════════════════════════════════════════════════════════════════════════════
%  ROUND 1
% ══════════════════════════════════════════════════════════════════════════════
\section{Warm-Up \& Project Understanding \textnormal{\textcolor{slatelight}{(5--8 min)}}}

\emph{These ease the candidate in and test whether they truly own the project vs.\ copied it.}

% ── Q1 ──
\subsection{Walk me through your project in 2 minutes. What problem does it solve, and who would use it?}

\begin{listenfor}
Can they articulate the purpose clearly? Do they mention MERN stack, auth, CRUD, dashboard? Or do they just say ``it's an expense tracker.''
\end{listenfor}

% ── Q2 ──
\subsection{Why did you choose the MERN stack (MongoDB, Express, React, Node) for this project? Would you make the same choice again?}

\begin{listenfor}
Any reasoning — ``I wanted to learn it'', ``JavaScript everywhere'', ``MongoDB is flexible for user-specific data''. Red flag if they can't explain \emph{why} MongoDB over, say, PostgreSQL.
\end{listenfor}

% ── Q3 ──
\subsection{What was the hardest part of building this project? Walk me through how you solved it.}

\begin{listenfor}
Genuine struggle stories — debugging CORS with cookies, getting Docker multi-stage builds to work, the currency conversion logic, pagination. A candidate who built this will have \emph{stories}.
\end{listenfor}

\newpage

% ══════════════════════════════════════════════════════════════════════════════
%  ROUND 2
% ══════════════════════════════════════════════════════════════════════════════
\section{Backend \& Database Deep-Dive \textnormal{\textcolor{slatelight}{(10--12 min)}}}

% ── Q4 ──
\subsection{Your \texttt{User} model stores passwords. Explain step-by-step what happens to a password from the moment a user types it in the registration form to when it's saved in MongoDB.}

\begin{listenfor}
Frontend sends it over HTTPS $\rightarrow$ Express receives the plain text $\rightarrow$ \texttt{express-validator} validates length $\rightarrow$ Mongoose \texttt{pre("save")} hook fires $\rightarrow$ \texttt{bcrypt.genSalt(12)} $\rightarrow$ \texttt{bcrypt.hash()} $\rightarrow$ hashed string stored. The password field has \texttt{select: false} so it's never returned in queries by default.
\end{listenfor}

\begin{followup}
``Why 12 salt rounds? What happens if you increase it to 20?''

\medskip
\emph{Expected:} Each round doubles computation time — exponential cost increase. 12 rounds is the current industry sweet spot balancing security and performance.
\end{followup}

% ── Q5 ──
\subsection{You're using JWTs stored in httpOnly cookies for authentication. Why httpOnly? Why not just store the token in localStorage?}

\begin{listenfor}
\texttt{httpOnly} cookies can't be accessed by JavaScript, which protects against XSS attacks stealing the token. \texttt{localStorage} is vulnerable to XSS.
\end{listenfor}

\begin{followup}
``Your code sets \texttt{sameSite: 'none'} in production but \texttt{'strict'} in development. Why the difference?''

\medskip
\emph{Expected:} Split deployments on different subdomains in production require \texttt{none}; local dev runs on the same origin so \texttt{strict} is safer.
\end{followup}

% ── Q6 ──
\subsection{In your Expense schema, you have compound indexes like \texttt{\{user: 1, date: -1\}} and \texttt{\{user: 1, category: 1, date: -1\}}. Why those specific combinations?}

\begin{listenfor}
Understanding that queries always filter by \texttt{user} first (data isolation), then sort by \texttt{date} descending (most recent first), and optionally filter by \texttt{category}. The compound index matches these query patterns exactly.
\end{listenfor}

\begin{followup}
``Your search feature uses \texttt{\{\$regex: search, \$options: 'i'\}}. Will your index help with that query?''

\medskip
\emph{Expected:} No — unanchored regex can't use B-tree indexes; it does a collection scan. For a personal tracker it's fine, but at scale you'd need a text index or Atlas Search.
\end{followup}

% ── Q7 ──
\subsection{Look at your \texttt{getBudget} controller — it fetches the budget \emph{and} runs an aggregation to calculate spending in a single API call. Why is that better than making two separate calls?}

\begin{listenfor}
Reduces network round-trips, keeps the UI simpler (one loading state), and the aggregation runs server-side in MongoDB which is faster than fetching raw data and computing in Node.js.
\end{listenfor}

% ── Q8 ──
\subsection{Your currency conversion endpoint (\texttt{PUT /api/auth/currency}) updates \emph{every} expense, income, budget, and goal in the database when a user switches currency. What could go wrong?}

\begin{listenfor}
\textbf{Atomicity} — if one \texttt{updateMany} fails midway, you end up with partially converted data. There's no transaction wrapping the multiple updates. Also, storing converted amounts permanently means you lose the original values. A better approach might be storing amounts in one canonical currency and converting at display time.
\end{listenfor}

\begin{followup}
``I see you used MongoDB's \texttt{\$mul} operator instead of loading all documents into Node.js. Why?''

\medskip
\emph{Expected:} Performance — \texttt{\$mul} runs directly in the database engine without transferring documents over the network.
\end{followup}

\newpage

% ══════════════════════════════════════════════════════════════════════════════
%  ROUND 3
% ══════════════════════════════════════════════════════════════════════════════
\section{Frontend \& React \textnormal{\textcolor{slatelight}{(8--10 min)}}}

% ── Q9 ──
\subsection{How does your frontend know if a user is logged in? Walk me through what happens when the app first loads.}

\begin{listenfor}
\texttt{AuthContext} provides user state. On mount, it calls \texttt{GET /api/auth/me}. If the httpOnly cookie is valid, the server returns user data; the context sets \texttt{user} and \texttt{isAuthenticated}. If not, the user is redirected to login. The cookie is sent automatically because \texttt{axios} is configured with \texttt{withCredentials: true}.
\end{listenfor}

% ── Q10 ──
\subsection{You have protected routes in your app. How do you prevent an unauthenticated user from accessing \texttt{/dashboard} or \texttt{/expenses}?}

\begin{listenfor}
A \texttt{ProtectedRoute} wrapper component that checks \texttt{isAuthenticated} from context and either renders \texttt{<Outlet />} or redirects to \texttt{/login}. Understanding that this is \emph{client-side} protection only — the real security is the \texttt{protect} middleware on the backend that validates the JWT.
\end{listenfor}

% ── Q11 ──
\subsection{You recently added Framer Motion animations. Your \texttt{Modal.tsx} now uses \texttt{<AnimatePresence>}. Why can't you just use CSS transitions for the same effect?}

\begin{listenfor}
CSS can handle enter animations, but \emph{exit} animations are hard — once a component is removed from the DOM (React unmounts it), CSS has no way to animate it ``out.'' \texttt{AnimatePresence} detects when children are removed and keeps them in the DOM long enough to run the exit animation before actually unmounting.
\end{listenfor}

\begin{followup}
``In your \texttt{AppShell.tsx}, you use \texttt{key=\{location.pathname\}} on \texttt{<motion.main>}. Why is the key important?''

\medskip
\emph{Expected:} React uses the key to determine identity. When the pathname changes, the key changes, React treats it as a \emph{new} component, unmounts the old one (triggering exit animation) and mounts the new one (triggering enter animation). Without the key, React would just re-render the same component with no animation.
\end{followup}

% ── Q12 ──
\subsection{Your Dashboard page makes multiple API calls (\texttt{/summary}, \texttt{/spending-by-category}, \texttt{/trend}, \texttt{/goals}). How do you handle the loading state? What does the user see while data loads?}

\begin{listenfor}
Skeleton loading components (\texttt{Skeleton}, \texttt{ActivityFeedSkeleton}). Each section loads independently — the hero balance can show up before the spending chart is ready. This is better than a single full-page spinner because it gives perceived performance.
\end{listenfor}

\newpage

% ══════════════════════════════════════════════════════════════════════════════
%  ROUND 4
% ══════════════════════════════════════════════════════════════════════════════
\section{API Design \& Error Handling \textnormal{\textcolor{slatelight}{(5--7 min)}}}

% ── Q13 ──
\subsection{Your \texttt{POST /api/expenses/:id/duplicate} endpoint clones an expense. Why did you make this a POST instead of, say, GET or PUT?}

\begin{listenfor}
POST because it \emph{creates} a new resource (side effect). GET should be idempotent and safe — calling it twice shouldn't create two copies. PUT is for updating an existing resource, not creating new ones. This shows understanding of REST semantics.
\end{listenfor}

% ── Q14 ──
\subsection{You use \texttt{express-validator} for input validation on create routes, but your update routes only rely on Mongoose's \texttt{runValidators: true}. Is that intentional? What's the difference?}

\begin{listenfor}
Honesty — this might be an oversight. \texttt{express-validator} runs \emph{before} hitting the database and gives cleaner error messages. Mongoose validators run at the database layer and throw different error shapes. Ideally both routes should have consistent validation.
\end{listenfor}

\begin{tipbox}
This question tests intellectual honesty. A great candidate will say ``That's actually inconsistent, I should fix that.'' A weaker one will try to justify it as intentional.
\end{tipbox}

% ── Q15 ──
\subsection{You have a global error handler middleware. What happens if an async controller throws an unhandled error? Show me the flow.}

\begin{listenfor}
\texttt{express-async-handler} wraps each controller $\rightarrow$ catches the rejected promise $\rightarrow$ passes the error to \texttt{next()} $\rightarrow$ Express routes it to the \texttt{errorHandler} middleware $\rightarrow$ which sends \texttt{\{success: false, message, stack\}} (stack hidden in production). Without \texttt{asyncHandler}, Express would crash or hang on unhandled promise rejections.
\end{listenfor}

\newpage

% ══════════════════════════════════════════════════════════════════════════════
%  ROUND 5
% ══════════════════════════════════════════════════════════════════════════════
\section{DevOps \& Deployment \textnormal{\textcolor{slatelight}{(5--7 min)}}}

% ── Q16 ──
\subsection{You have a two-stage Dockerfile for the backend. Explain what multi-stage builds are and why you use them.}

\begin{listenfor}
Stage 1 installs all dependencies (including devDependencies like TypeScript) and compiles. Stage 2 starts fresh with only production dependencies and the compiled output. The final image is smaller (no TS compiler, no dev tools), more secure (less attack surface), and faster to deploy.
\end{listenfor}

% ── Q17 ──
\subsection{Your \texttt{docker-compose.prod.yml} adds a Caddy reverse proxy. Why not just expose the Nginx frontend directly on port 443?}

\begin{listenfor}
Caddy provides \emph{automatic} HTTPS — it handles Let's Encrypt certificate provisioning, renewal, and HTTP$\rightarrow$HTTPS redirect with zero configuration. Doing this manually with Nginx requires certbot, cron jobs, config files. Also, the two-tier architecture (Caddy $\rightarrow$ Nginx $\rightarrow$ Express) separates concerns: Caddy handles TLS termination, Nginx handles static serving and API proxying.
\end{listenfor}

% ── Q18 ──
\subsection{In your frontend Dockerfile, you deliberately skip copying \texttt{package-lock.json}. That's unusual — why?}

\begin{listenfor}
There's a documented reason — Windows-generated lockfiles sometimes miss Linux-specific native binaries like \texttt{@rollup/rollup-linux-x64-musl} (npm issue \#4828). Running \texttt{npm install} from scratch inside the Alpine container ensures the correct platform-specific packages are resolved. This is a real-world cross-platform gotcha.
\end{listenfor}

\newpage

% ══════════════════════════════════════════════════════════════════════════════
%  ROUND 6
% ══════════════════════════════════════════════════════════════════════════════
\section{Thinking \& Growth \textnormal{\textcolor{slatelight}{(5 min)}}}

% ── Q19 ──
\subsection{If this app had 10,000 users and each had 5,000 expenses, what's the first thing that would break?}

\begin{listenfor}
Thoughtful reasoning, not necessarily the ``right'' answer. Good answers include:
\begin{itemize}[nosep]
  \item The regex search doing collection scans
  \item The dashboard aggregation queries getting slow without proper indexing
  \item The in-place currency conversion updating millions of documents
  \item MongoDB Atlas free tier limits
  \item No caching layer (Redis) for repeated dashboard calls
  \item Single Node.js process (no clustering / PM2)
\end{itemize}
\end{listenfor}

% ── Q20 ──
\subsection{If I gave you one more week on this project, what would you build next and why?}

\begin{listenfor}
Self-awareness about gaps in their own codebase. Strong answers might reference:
\begin{itemize}[nosep]
  \item Actually sending the password reset emails (currently a TODO)
  \item Wiring up the \texttt{multer} and \texttt{node-cron} dependencies that are installed but unused
  \item Adding CI/CD (GitHub Actions)
  \item Writing tests (the test infrastructure is set up but no tests exist)
  \item Adding a root \texttt{.gitignore} to prevent accidental \texttt{.env} commits
  \item Building out the Goals page (currently a stub despite the API being complete)
  \item Removing unused dependencies (\texttt{chart.js}, \texttt{react-chartjs-2})
  \item Budget notification alerts that actually trigger
\end{itemize}
\end{listenfor}

\newpage

% ══════════════════════════════════════════════════════════════════════════════
%  SCORING RUBRIC
% ══════════════════════════════════════════════════════════════════════════════
\section*{Scoring Rubric}
\addcontentsline{toc}{section}{Scoring Rubric}

{\small\textcolor{slatelight}{\emph{Internal — not shared with candidate.}}}

\vspace{0.5cm}

\renewcommand{\arraystretch}{1.6}
\begin{tabularx}{\textwidth}{c X}
  \toprule
  \textbf{Score} & \textbf{Description} \\
  \midrule
  $\bigstar\bigstar\bigstar\bigstar\bigstar$
    & Answers most questions confidently. Can explain \emph{why} behind decisions.
      Acknowledges what they'd improve. Thinks about trade-offs. \\
  $\bigstar\bigstar\bigstar\bigstar\phantom{\bigstar}$
    & Solid understanding of 70\%+ of their codebase. Gets tripped up on deeper
      questions (indexing, atomicity) but reasons through them with hints. \\
  $\bigstar\bigstar\bigstar\phantom{\bigstar\bigstar}$
    & Knows what the code does but not always \emph{why}. Struggles with security
      and performance questions. Likely followed tutorials closely. \\
  $\bigstar\bigstar\phantom{\bigstar\bigstar\bigstar}$
    & Can describe the project at a surface level but can't explain implementation
      details. Likely copy-pasted significant portions. \\
  $\bigstar\phantom{\bigstar\bigstar\bigstar\bigstar}$
    & Cannot explain basic project decisions. Red flag for authorship. \\
  \bottomrule
\end{tabularx}

\vspace{1cm}

\begin{tipbox}
\textbf{For a fresher}, scoring $\bigstar\bigstar\bigstar$--$\bigstar\bigstar\bigstar\bigstar$ is excellent. The goal isn't perfection — it's to see genuine ownership, curiosity, and the ability to reason about code they wrote. A candidate who says ``I'm not sure, but I think it might be because\ldots'' is far better than one who freezes or guesses confidently with wrong answers.
\end{tipbox}

\end{document}
